\chapter{Équations différentielles et équations intégro-différentielles}

Un ensemble très important des problèmes scientifiques est décrit par des
équations différentielles ordinaires et des équations aux dérivées partielles
linéaires et non linéaires, à conditions aux limites (ou aux frontières) ou à
conditions initiales. Dans certains cas, un tel problème peut se résoudre en
trouvant les solutions générales et particulières par les méthodes analytiques
et classiques (résolution directe, utilisation des facteurs intégrants,
utilisation de la transformée de Laplace, etc.).

Lorsqu'aucune méthode analytique ne permet de calculer la solution analytique,
il faut trouver les méthodes numériques pour approcher la valeur cherchée.

Ce chapitre est consacré à la présentation des méthodes numériques pour
résoudre les équations différentielles ordinaires à valeurs initiales et à
valeurs aux limites. On abordera aussi les méthodes de résolution numérique
des équations intégro-différentielles dans lesquelles il y a la présence des
dérivées de la variable inconnue et des intégrales de la variable inconnue.

\section{Généralités}

\subsection{Définition et classification}

Une équation différentielle est une relation qui lie une fonction $y$ à ses
dérivées. S'il y a une seule variable indépendante, on parle alors d'équation
différentielle ordinaire. Dans le cas de plusieurs variables indépendantes, on
parle d'équations aux dérivées partielles. Les équations différentielles
ordinaires de variable indépendante $x$ peuvent se mettre sous la forme
\begin{equation}
  F\bigl(x, y, y', \dots, y^{(n-1)}, y^{(n)}\bigr) = 0
  \tag{1}
\end{equation}
où $y^{(i)}$ est la dérivée d'ordre $i$ et $n$ est l'ordre de l'équation
différentielle et $a \le x \le b$.

Le degré d'une équation différentielle, si celle-ci peut être décrite comme un
polynôme, où les indéterminées sont des dérivées, est le degré de la dérivée de
l'ordre le plus élevé. Par exemple, l'équation
$(y'')^3 + (y')^5 + 7y = e^x$ est une équation différentielle de second ordre
et de degré 3. Si la fonction $F$ est un polynôme dont les degrés de $y$ et de
ses dérivées sont égaux à 0 ou 1 et si $F$ ne contient pas des produits de $y$
et de ses dérivées, ou des dérivées entre elles, alors l'équation
différentielle est dite \emph{linéaire}. Dans tous les autres cas, l'équation
est dite \emph{non linéaire}.

La résolution d'une équation différentielle d'ordre $n$ nécessite la
connaissance au préalable de $n$ constantes, notamment $n$ conditions imposées
à $y$ et à ses dérivées. Si l'on fixe la valeur de $y$ et de ses $(n-1)$
dérivées au point initial $x = a$, alors, on a un \emph{problème aux valeurs
initiales}. Si par contre les $n$ conditions sont imposées les unes en $x = a$
et les autres en $x = b$, on a un \emph{problème aux valeurs aux limites}.

\subsection{Origine des équations différentielles}

Les équations différentielles sont d'une importance fondamentale dans presque
tous les domaines de la science. Ceci est dû au fait que plusieurs lois et
phénomènes ont une formulation mathématique sous forme d'équations
différentielles. Par exemple, la seconde loi de Newton est une équation
différentielle du second ordre. De même la loi de la radioactivité est une
équation du premier ordre.

\subsection{Position du problème numérique}

Considérons l'équation différentielle ordinaire du premier ordre
\begin{equation}
  y' = f(x, y) .
  \tag{2}
\end{equation}
Le problème est de trouver la valeur $\bar{y}$ de $y$ correspondant à
$x = x_0 + h$ (ou en $x_k = x_0 + kh$) connaissant $y = y_0$ pour $x = x_0$.
$h$ est un nombre réel et $k$ un nombre entier.

En intégrant (2), on obtient
\begin{equation}
  y = y_0 + \int_{x_0}^{x} f(x, y)\,\dd x .
  \tag{3}
\end{equation}
On a alors
\begin{equation}
  \bar{y} = y_0 + \int_{x_0}^{x_0+h} f(x, y)\,\dd x = y_0 + L
  \qquad\text{où}\qquad
  L = \int_{x_0}^{x_0+h} f(x, y)\,\dd x .
  \tag{4}
\end{equation}
On peut donc utiliser des méthodes numériques de calcul des intégrales pour
calculer l'intégrale. On peut également utiliser le développement de
$y(x+h)$ en série de puissances de $h$ (puis insérer les différentes dérivées
à partir de l'équation (2)).

\section{Méthodes pour les problèmes à valeurs initiales}

\subsection{Méthode d'Euler}

Encore appelée méthode de la dérivée première, elle s'exprime par la formule
\[
  y(x+h) = y(x) + hf(x, y) .
\]
La formule itérative est donc
\[
  y_{k+1} = y_k + hf(x_k, y_k) .
\]
C'est la méthode la plus simple. Mais son inconvénient réside dans le fait
qu'elle est souvent très instable et exige des pas de calcul $h$ très petits.

\paragraph{Algorithme de la méthode d'Euler}

\begin{algo}
  \State Définir $f(x,y)$
  \State $x \leftarrow x_0$, $y \leftarrow y_0$
  \State Donner la valeur de $h$
  \For{$k$ allant de $1$ à $n$}
    \State $y \leftarrow y + hf(x,y)$
    \State $x \leftarrow x + h$
    \Print $x$ et $y$
  \EndFor
\end{algo}

\subsection{Les méthodes de Runge et de Kutta}

Ce sont les méthodes élaborées par C.~Runge en 1894 et améliorées par
W.~Kutta vers 1901. Elles utilisent les méthodes d'intégration numérique des
trapèzes et de Simpson. Ces méthodes sont les plus utilisées car elles sont
très stables.

\subsubsection{Méthode de Runge-Kutta du second ordre}

Elle est basée sur la formule intégrale des trapèzes (voir chapitre sur
l'intégration numérique) et son schéma itératif se met sous la forme :
\[
  y_{k+1} = y_k + \frac{h}{2}\Bigl\{f(x_k, y_k) + f\bigl[x_k + h,\; y_k + hf(x_k, y_k)\bigr]\Bigr\} .
\]
Elle présente une erreur d'ordre $O(h^3)$.

\begin{exercice}
Écrire l'algorithme de la méthode de Runge-Kutta d'ordre 2.
\end{exercice}

\subsubsection{Méthode de Runge du quatrième ordre}

Supposons les valeurs $y_0$, $y_1$, $y_2$ de $y$ connues aux points $x_0$,
$x_1 = x_0 + \dfrac{h}{2}$ et $x_2 = x_0 + h$.

En utilisant la formule d'intégration numérique de Simpson (voir chapitre
précédent), on obtient
\[
  L = \frac{h}{6}\Bigl[f(x_0, y_0) + 4f\Bigl(x_0 + \frac{h}{2}, y_1\Bigr) + f(x_0 + h, y_2)\Bigr].
\]
Les valeurs $y_1$ et $y_2$ étant inconnues, on utilise les approximations
\[
  y_1 \approx y_0 + \frac{h}{2}f(x_0, y_0) = y_0 + \frac{h}{2}f_0, \qquad
  y_2 \approx y_0 + hf(x_0 + h, y_0 + hf_0) .
\]
Ces approximations sont choisies de telle manière qu'en développant $L$ en
séries de puissance de $h$, ces 3 premiers termes coïncident avec le
développement en série de Taylor de $y(x+h)$. Il vient alors
\[
  L \approx \frac{h}{6}\Bigl\{f(x_0, y_0)
     + 4f\Bigl[x_0 + \frac{h}{2},\; y_0 + \frac{h}{2}f(x_0, y_0)\Bigr]
     + f\Bigl[x_0 + h,\; y_0 + hf\bigl(x_0 + h, y_0 + hf(x_0, y_0)\bigr)\Bigr]\Bigr\}.
\]
Soit plus généralement,
\[
  y_{k+1} = y_k + L
\]
avec
\[
  L = \frac{h}{6}\Bigl\{f(x_k, y_k)
     + 4f\Bigl[x_k + \frac{h}{2},\; y_k + \frac{h}{2}f(x_k, y_k)\Bigr]
     + f\Bigl[x_k + h,\; y_k + hf\bigl(x_k + h, y_k + hf(x_k, y_k)\bigr)\Bigr]\Bigr\}.
\]
Elle présente une erreur d'ordre $O(h^5)$.

Pratiquement on fait le calcul de la manière suivante. On calcule les quantités
$L_i$ définies par :
\[
  \begin{aligned}
    L_1 &= hf(x_k, y_k)\\
    L_2 &= hf(x_k + h, y_k + L_1)\\
    L_3 &= hf(x_k + h, y_k + L_2)\\
    L_4 &= hf\Bigl(x_k + \frac{h}{2}, y_k + \frac{L_1}{2}\Bigr)
  \end{aligned}
\]
Par la suite, on écrit
\[
  y_{k+1} = y_k + \frac{1}{6}\bigl(L_1 + 4L_4 + L_3\bigr).
\]

\begin{exercice}
Établir l'algorithme de la méthode de Runge d'ordre 4.
\end{exercice}

\subsubsection{Méthode de Runge-Kutta ou de Kutta-Simpson du quatrième ordre}

C'est une amélioration de la méthode de Runge d'ordre 4. Elle présente
également une erreur d'ordre $O(h^5)$ et son schéma itératif s'écrit :
\[
  y_{k+1} = y_k + \frac{1}{6}\bigl(L_1 + 2L_2 + 2L_3 + L_4\bigr)
\]
avec
\[
  \begin{aligned}
    L_1 &= hf(x_k, y_k)\\
    L_2 &= hf\Bigl(x_k + \frac{h}{2}, y_k + \frac{L_1}{2}\Bigr)\\
    L_3 &= hf\Bigl(x_k + \frac{h}{2}, y_k + \frac{L_2}{2}\Bigr)\\
    L_4 &= hf(x_k + h, y_k + L_3)
  \end{aligned}
\]
C'est la méthode la plus utilisée. Il existe d'autres variantes des méthodes de
Runge-Kutta telles que la méthode de Kutta-Merson ou de Kutta-Simpson du
sixième ordre.

\subsection{Méthodes à pas multiples ou d'Adams-Bashforth}

\subsubsection{Formule du second ordre}

Soit l'équation différentielle
\[
  y' = f(x, y) .
\]
On veut connaître $y(x+h)$ connaissant $y(x)$ et $y(x-h)$. Pour cela, on fait
le développement suivant :
\[
  y(x+h) = y(x) + hy' + \frac{h^2}{2}y'' + O(h^3)
  \quad\Longrightarrow\quad
  y_{k+1} = y_k + hy'_k + \frac{h^2}{2}y''_k .
\]
On a aussi
\[
  y'(x-h) = y'(x) - hy''(x) .
\]
D'où, il vient alors
\[
  y''_k = \frac{y'_k - y'_{k-1}}{h} .
\]
Finalement, on obtient
\[
  y_{k+1} = y_k + \frac{h}{2}\bigl[3y'_k - y'_{k-1}\bigr]
  \quad\Longrightarrow\quad
  y_{k+1} = y_k + \frac{h}{2}\bigl[3f(x_k, y_k) - f(x_{k-1}, y_{k-1})\bigr].
\]
C'est la formule d'Adams-Bashforth du second ordre présentant une erreur
d'ordre $O(h^3)$.

Pratiquement au début des itérations, il faut connaître les valeurs de $y$ en
$x_0$ et en $x_1$ afin de calculer la première valeur $y_2$ de l'itération.
Pour avoir $y_1$, on peut utiliser l'une des formules approchées classiques.
Par exemple en utilisant la formule d'intégration des trapèzes, on obtient :
\[
  y_1 = y_0 + \frac{h}{2}\Bigl[f(x_0, y_0) + f\bigl[x_0 + h, y_0 + hf(x_0, y_0)\bigr]\Bigr].
\]
La formule itérative peut alors débuter par $k = 1$, et on calcule les $y_2$,
$y_3$, $y_4$\dots{} Par exemple,
\[
  y_2 = y_1 + \frac{h}{2}\bigl[3f(x_1, y_1) - f(x_0, y_0)\bigr].
\]

\begin{exercice}
Écrire l'algorithme de la méthode d'Adams-Bashforth du second ordre.
\end{exercice}

\subsubsection{Formules d'ordre supérieur}

Pour obtenir des formules d'ordre supérieur présentant par exemple des erreurs
d'ordre $O(h^4)$ ou $O(h^5)$, il suffit de développer à l'ordre supérieur
$y(x+h)$ et de remplacer les différentes dérivées.

\subsection{Méthode de prédicteur-correcteur}

Les méthodes de Runge, Runge-Kutta et de Kutta-Simpson sont fondamentalement
implicites car l'inconnue $y_{k+1}$ est exprimée à l'aide d'une fonction
contenant $y_{k+1}$. Par exemple, dans le cas de Runge-Kutta d'ordre 2, on a
normalement
\[
  y_{k+1} = y_k + \frac{h}{2}\bigl[f(x_k, y_k) + f(x_{k+1}, y_{k+1})\bigr].
\]
Afin d'avoir une expression explicite, on part des approximations pour
exprimer $y_{k+1}$ du second membre de l'égalité en fonction de $y_k$. On
pourrait aussi partir des formules d'Adams-Bashforth pour évaluer (prédire) les
valeurs de $y_{k+1}$, puis insérer ces valeurs prédites dans le second membre
de la formule implicite afin d'avoir une valeur corrigée. Dans un tel
processus, la formule explicite est appelée \emph{prédicteur} tandis que la
formule implicite utilisée est appelée \emph{correcteur}. Par exemple une
méthode de prédicteur-correcteur peut utiliser les formules suivantes :
\begin{align*}
  y_{k+1} &= y_k + \frac{h}{2}\bigl[3f(x_k, y_k) - f(x_{k-1}, y_{k-1})\bigr]\\
          &\equiv \text{prédicteur (formule d'Adams-Bashforth du second ordre),}\\[0.6em]
  y_{k+1} &= y_k + \frac{h}{2}\bigl[f(x_k, y_k) + f(x_{k+1}, y_{k+1})\bigr]\\
          &\equiv \text{correcteur (Runge-Kutta d'ordre 2).}
\end{align*}
On peut également combiner les autres formules implicites de RK aux autres
formules explicites d'Adams-Bashforth pour avoir d'autres méthodes de
prédicteur-correcteur plus élaborées.

\begin{exercice}
Établir l'algorithme de la méthode de prédicteur-correcteur, utilisant comme
prédicteur la formule d'Adams-Bashforth du second ordre et comme correcteur la
formule de $RK_2$ et traduire cet algorithme en Pascal ou en Fortran ou en C.
\end{exercice}

\section{Méthodes pour les problèmes aux valeurs limites}

Il s'agit des problèmes de la forme
\[
  y^{(n)} = f\bigl(x, y, y', \dots, y^{(n-1)}\bigr), \qquad a < x < b
\]
avec $y(a) = y_1$ et $y(b) = y_2$ et d'autres conditions sur les dérivées. On
peut résoudre ce type de problème par plusieurs méthodes dont la méthode de
tir, la méthode des fonctions complémentaires ou de superposition et la
méthode des différences finies.

\subsection{Méthode de tir}

Elle est utilisée dans le cas des problèmes aux limites non linéaires. Il
s'agit de transformer un problème aux valeurs aux limites en un problème aux
valeurs initiales par un choix adéquat des valeurs initiales. Ce choix se fait
à l'une des bornes du domaine d'intégration et doit être tel que la solution
passe le plus près possible de la valeur de la grandeur à l'autre borne du
domaine d'intégration.

Par exemple, considérons une équation du second ordre
\[
  y'' = f(x, y, y') \quad\text{avec}\quad y(a) = \eta_1 \quad\text{et}\quad y(b) = \eta_2 .
\]
La méthode de tir consiste à choisir $\gamma$ telle que $y'(a) = \gamma$ et de
résoudre le problème aux valeurs initiales
\[
  y'' = f(x, y, y') \quad\text{avec}\quad y(a) = \eta_1 \quad\text{et}\quad y'(a) = \gamma .
\]
Le choix de $\gamma$ est bon lorsqu'en $x = b$, on a $y(b) \approx \eta_2$.
Sinon il faut choisir une autre valeur pour $\gamma$. Une procédure à utiliser
pour le choix de $\gamma$ est la suivante. On sait que la valeur $y(b)$ dépend
de $\gamma$ ; c'est-à-dire que $y(b) = y(b, \gamma)$. Pour un bon choix de
$\gamma$, on doit avoir
\[
  y(b, \gamma) - \eta_2 \approx 0 \;\Longleftrightarrow\; f(\gamma) = 0
  \quad\text{où}\quad f(\gamma) = y(b, \gamma) - \eta_2 .
\]
Or d'après la méthode de Newton (voir chapitre II), le schéma itératif donnant
la solution de cette équation algébrique est
\[
  \gamma_{n+1} = \gamma_n - \frac{f(\gamma_n)}{f'(\gamma_n)}
  = \gamma_n - \frac{y(b, \gamma_n) - \eta_2}{y'(b, \gamma_n)}
  = \gamma_n - \frac{(\gamma_n - \gamma_{n-1})\bigl(y(b, \gamma_n) - \eta_2\bigr)}
                    {y(b, \gamma_n) - y(b, \gamma_{n-1})}
\]

Ainsi en pratique, on commence par donner une valeur initiale $\gamma_0$ à
$\gamma$. Ensuite, on choisit la valeur suivante $\gamma_1$, les autres valeurs
de $\gamma$ $(\gamma_2, \gamma_3, \dots)$ sont ensuite évaluées à partir de la
formule ci-dessus.

\begin{exercice}
Établir l'algorithme de la méthode de tir.
\end{exercice}

\subsection{Méthode des fonctions complémentaires ou de superposition}

C'est une méthode qui transforme une équation différentielle linéaire à
valeurs aux limites en deux problèmes différentiels à valeurs initiales. Elle a
plusieurs variantes que nous exposerons sur deux équations différentielles du
second ordre.

\subsubsection{Première variante}

Considérons l'équation
\[
  \left\{
  \begin{aligned}
    &y'' = p(x)y' + q(x)y + r(x)\\
    &y(a) = \eta_1 \quad\text{et}\quad y(b) = \eta_2
  \end{aligned}
  \right.
\]
Supposons que :
\[
  y(x) = y_1(x) + \mu y_2(x) \qquad (\mu = \text{cte} \neq 0).
\]
Il vient
\[
  \begin{aligned}
    y''_1 &= p(x)y'_1 + q(x)y_1 + r(x)\\
    y''_2 &= p(x)y'_2 + q(x)y_2
  \end{aligned}
\]
En $x = a$, on a $y_1(a) + \mu y_2(a) = \eta_1$.

Imposons que $y_1(a) = \eta_1 \Longrightarrow y_2(a) = 0$.

En $x = b$, on a $y_1(b) + \mu y_2(b) = \eta_2$. D'où
\[
  \Longrightarrow\quad \mu = \frac{\eta_2 - y_1(b)}{y_2(b)} .
\]
Fixons
\[
  y'_1(a) = 0 \qquad\text{et}\qquad y'_2(a) = 1 .
\]
On obtient alors les deux problèmes différentiels à valeurs initiales
suivants :
\[
  \left\{
  \begin{aligned}
    &\left\{
    \begin{aligned}
      &y''_1 = p(x)y'_1 + q(x)y_1 + r(x)\\
      &y_1(a) = \eta_1, \quad y'_1(a) = 0
    \end{aligned}
    \right.\\[0.4em]
    &\left\{
    \begin{aligned}
      &y''_2 = p(x)y'_2 + q(x)y_2\\
      &y_2(a) = 0, \quad y'_2(a) = 1
    \end{aligned}
    \right.
  \end{aligned}
  \right.
  \qquad x \in\, ]a, b]
\]
On résout alors numériquement ces deux problèmes à valeurs initiales. On
obtient les valeurs de $y_1(b)$ et $y_2(b)$ et par conséquent la valeur de
$\mu$ et les différentes valeurs de $y$ aux n\oe uds de discrétisation
(puisque les valeurs de $y_1$ et $y_2$ sont connues en tout point de
l'intervalle $]a,b]$).

\subsubsection{Deuxième variante}

Considérons le problème suivant
\begin{align}
  y'' &= p(x)y + q(x) \tag{a}\\
  y'(a) &= \alpha_{00}y(a) + \alpha_{10} \tag{b}\\
  y'(b) &= \beta_{00}y(b) + \beta_{10} \tag{c}
\end{align}
Considérons également l'équation
\begin{equation}
  y' = \alpha_0(x)y + \alpha_1(x) \tag{d}
\end{equation}
où nous avons choisi $\alpha_0(x)$ et $\alpha_1(x)$ tels que $y$ satisfait
toujours l'équation (a). En différenciant l'équation (d) par rapport à $x$ et
en identifiant à l'équation (a), on obtient
\begin{align}
  \alpha'_0 + \alpha_0^2 &= p(x) \tag{e}\\
  \alpha'_1 + \alpha_0\alpha_1 &= q(x) \tag{f}
\end{align}
avec $\alpha_0(a) = \alpha_{00}$, $\alpha_1(a) = \alpha_{10}$.

Les équations (e) et (f) sont des équations différentielles du
1\textsuperscript{er} ordre que l'on peut résoudre analytiquement ou
numériquement pour trouver les valeurs de $\alpha_0(x)$ et de $\alpha_1(x)$.

De l'équation (d), on a
$y'(b) = \alpha_0(b)y(b) + \alpha_1(b) = \beta_{00}y(b) + \beta_{10}$. On
obtient alors
\begin{align}
  y(b) &= \frac{\beta_{10} - \alpha_1(b)}{\alpha_0(b) - \beta_{00}} \tag{g}\\
  y'(b) &= \frac{\beta_{00}\alpha_1(b) - \beta_{10}\alpha_0(b)}{\beta_{00} - \alpha_0(b)} \tag{h}
\end{align}
L'équation (a) est ainsi transformée en un problème à valeurs initiales. On
peut en effet résoudre le problème $y'' = p(x)y + q(x)$ avec les conditions
initiales (g) et (h) ou résoudre le problème $y' = \alpha_0(x)y + \alpha_1(x)$
avec la condition initiale (g).

\subsection{Méthode des différences finies}

Elle convertit le problème de résolution d'une équation différentielle à la
résolution des équations algébriques.

Soit le problème :
\[
  y'' = p(x)y' + q(x)y + r(x) \qquad\text{avec}\quad y(a) = \eta_1 \quad\text{et}\quad y(b) = \eta_2 .
\]
On a :
\[
  y'(x) = \frac{y(x+h) - y(x)}{h}
  \qquad\text{et}\qquad
  y''(x) = \frac{y(x+h) - 2y(x) + y(x-h)}{h^2} .
\]
Soit $x = x_k$, on posera
\[
  y(x_k) = y_k, \quad p(x_k) = p_k, \quad q(x_k) = q_k, \quad r(x_k) = r_k .
\]
Par substitution des expressions des dérivées dans l'équation différentielle,
on aboutit à l'équation discrète
\[
  (1 - hp_k)\,y_{k+1} - \bigl(2 - hp_k + h^2q_k\bigr)\,y_k + y_{k-1} = h^2r_k
\]
pour $k = 1$ à $N-1$ et $h = \dfrac{b - a}{N}$.

On aboutit ainsi à un système d'équations algébriques linéaires de la forme
\[
  \left\{
  \begin{aligned}
    &y_0 = \eta_1\\
    &(1 - hp_1)\,y_2 - \bigl(2 - hp_1 + h^2q_1\bigr)\,y_1 + y_0 = h^2r_1\\
    &(1 - hp_2)\,y_3 - \bigl(2 - hp_2 + h^2q_2\bigr)\,y_2 + y_1 = h^2r_2\\
    &\qquad\vdots\\
    &(1 - hp_{N-1})\,y_N - \bigl(2 - hp_{N-1} + h^2q_{N-1}\bigr)\,y_{N-1} + y_{N-2} = h^2r_{N-1}\\
    &y_N = \eta_2
  \end{aligned}
  \right.
\]
Ce système algébrique peut être résolu par les méthodes itératives de Jacobi
ou de Gauss-Seidel.

\section{Les équations intégrales et intégro-différentielles}

\subsection{Définition}

On appelle équation intégrale une équation dont la fonction inconnue se trouve
dans le symbole d'intégration.

C'est l'exemple de l'équation suivante :
\[
  y(x) = e^x\sin x + 2\int_0^x \cos(x - t)\,y(t)\,\dd t .
\]
On a aussi les équations intégro-différentielles comme par exemple
\[
  y'(x) = 3y(x) + 2x + \int_0^x \sin w(x - t)\,y(t)\,\dd t .
\]
Une équation intégrale est dite linéaire lorsque la fonction inconnue $y$
intervient avec un degré 1. Sous la forme générale, les équations intégrales
linéaires se présentent de la manière suivante
\[
  y(x) = f(x) + \int_0^x k(x,t)\,y(t)\,\dd t .
\]
La fonction $k(x,t)$ est appelée \emph{noyau} de l'équation.

Dans les cas simples, les équations intégrales peuvent se résoudre
analytiquement en utilisant par exemple la méthode de la fonction de Gauss, de
la transformée de Laplace, de la transformation de Mellin, ou par les
approximations successives. Dans le cas contraire, il faut utiliser les
méthodes numériques.

\subsection{Résolution numérique des équations intégrales linéaires}

Soit l'équation linéaire
\begin{equation}
  y(x) = f(x) + \int_a^b k(x,t)\,y(t)\,\dd t .
  \tag{1}
\end{equation}
L'intégrale définie dans (1) peut être remplacée par les formules des trapèzes
et de Simpson. On obtient alors
\begin{equation}
  y(x) = f(x) + (b - a)\bigl[c_1k(x,t_1)y(t_1) + c_2k(x,t_2)y(t_2) + \dots
         + c_nk(x,t_n)y(t_n)\bigr]
  \tag{2}
\end{equation}
où $t_1, t_2, \dots, t_n$ sont les n\oe uds de subdivision de l'intervalle
$[a,b]$ et les coefficients $c_i$ dépendent de la formule d'intégration
utilisée.

Puisque l'équation (2) est vérifiée pour tout $x \in [a,b]$, elle est également
vérifiée pour $x = t_1$, $x = t_2$, \dots, $x = t_n$. Par conséquent, à partir
de l'équation (2), on obtient un système d'équations de la forme
\begin{equation}
  y(t_i) = f(t_i) + (b - a)\bigl[c_1k(t_i,t_1)y(t_1) + c_2k(t_i,t_2)y(t_2) + \dots
           + c_nk(t_i,t_n)y(t_n)\bigr]
  \tag{3}
\end{equation}
pour $i$ allant de $1$ à $n$ ($t_1 = a$, \dots, $t_i = a + (i-1)h$, \dots,
$t_n = b$).

Posons
\[
  y(t_i) = y_i \qquad\text{et}\qquad f(t_i) = f_i .
\]
Le système (3) devient alors
\begin{equation}
  \left\{
  \begin{aligned}
    y_1 &= f_1 + (b - a)\bigl[c_1k(t_1,t_1)y_1 + c_2k(t_1,t_2)y_2 + \dots + c_nk(t_1,t_n)y_n\bigr]\\
    y_2 &= f_2 + (b - a)\bigl[c_1k(t_2,t_1)y_1 + c_2k(t_2,t_2)y_2 + \dots + c_nk(t_2,t_n)y_n\bigr]\\
    &\;\;\vdots\\
    y_n &= f_n + (b - a)\bigl[c_1k(t_n,t_1)y_1 + c_2k(t_n,t_2)y_2 + \dots + c_nk(t_n,t_n)y_n\bigr]
  \end{aligned}
  \right.
  \tag{4}
\end{equation}
C'est un système de $n$ équations à $n$ inconnues que sont les $y_i$. La
résolution de ce système permet de déterminer les $y_i$ et en les substituant
dans l'équation (2), on obtient l'expression analytique de $y(x)$, solution de
l'équation (1).

\paragraph{Remarque} On peut transformer une équation différentielle
ordinaire en une équation intégrale. En résolvant l'équation intégrale, on
trouve alors une solution analytique de l'équation différentielle.

\paragraph{Exemple} Soit l'équation
\[
  y(x) = \frac{5x}{6} - \frac{1}{9} + \frac{1}{3}\int_0^1 (t + x)\,y(t)\,\dd t .
\]
Ici, on a $f(x) = 5x/6 - 1/9$ et $k(x,t) = x + t$.

Pour résoudre numériquement cette équation, utilisons la méthode de Simpson et
considérons pour cela trois points : $t_1 = 0$, $t_2 = 0{,}5$ et $t_3 = 1$. Le
système algébrique à résoudre est alors défini par :
\[
  y(x) = \frac{5x}{6} - \frac{1}{9}
       + \frac{1}{3}\times\frac{1}{3}\times\frac{1}{2}
         \Bigl[(t_1 + x)y(t_1) + 4(t_2 + x)y(t_2) + (t_3 + x)y(t_3)\Bigr]
\]
car $h = 1/2$.

En remplaçant $x$ par $t_i$, on obtient
\[
  \left\{
  \begin{aligned}
    y_1 &= -\frac{1}{9} + \frac{1}{18}\bigl(2y_2 + y_3\bigr) && \text{en } t = 0\\
    y_2 &= \frac{5}{12} - \frac{1}{9} + \frac{1}{18}\Bigl(\frac{y_1}{2} + 4y_2 + \frac{3}{2}y_3\Bigr) && \text{en } t = 0{,}5\\
    y_3 &= \frac{5}{6} - \frac{1}{9} + \frac{1}{18}\bigl(y_1 + 6y_2 + 2y_3\bigr) && \text{en } t = 1
  \end{aligned}
  \right.
\]
On trouve $y_1 = 0$, $y_2 = \frac{1}{2}$ et $y_3 = 1$. En remplaçant
$y_1 = y(t_1)$, $y_2 = y(t_2)$ et $y_3 = y(t_3)$ dans l'expression de $y(x)$,
on obtient $y(x) = x$.

\begin{exercice}
Discrétiser l'équation intégro-différentielle suivante :
\[
  y'(x) = 2y(x) + f(x) + \int_0^x k(x,t)\,y(t)\,\dd t .
\]
\end{exercice}
